\chapter{Fournisseurs de métadonnées}
\label{ch:providers}

Lorsque vous scannez un ISBN ou un EAN-13, mybibli ne cherche pas le
titre dans un catalogue local — il interroge un ou plusieurs
\emph{fournisseurs de métadonnées}\index{fournisseurs de métadonnées}
externes jusqu'à ce que l'un d'eux renvoie une réponse propre. Ce
chapitre explique la chaîne, où récupérer les clés d'API et comment
les faire tourner en toute sécurité.

\section{La chaîne de fournisseurs}

mybibli parcourt les fournisseurs dans l'ordre le mieux adapté au
type de média deviné :

\begin{longtable}{@{} l l p{4.5cm} c @{}}
\toprule
\textbf{Fournisseur} & \textbf{Types de média} & \textbf{Notes} & \textbf{Clé ?} \\
\midrule
\endhead
BnF\index{BnF}                  & Livres         & Catalogue de la
Bibliothèque nationale de France ; idéal pour les titres en français. & non \\
Library of Congress\index{Library of Congress} & Livres & Catalogue
de la bibliothèque nationale des États-Unis ; pendant anglophone de
la BnF. Fournit aussi les zones MARC~21. & non \\
Open Library\index{Open Library} & Livres         & Catalogue
bénévole ; large couverture anglaise / multilingue. & non \\
Google Books\index{Google Books} & Livres         & Commercial ; bon
sur les essais grand public. & optionnelle \\
MusicBrainz\index{MusicBrainz}  & Audio          & Métadonnées musicales
bénévoles ; très complet sur CD / vinyles. & non \\
OMDb\index{OMDb}                & Films, séries  & Palier gratuit ; quotas. & requise \\
TMDb\index{TMDb}                & Films, séries  & Gratuit avec attribution. & requise \\
BDGest\index{BDGest}            & BD, comics     & Catalogue BD
francophone (web-scraped). & non \\
\bottomrule
\end{longtable}

Le premier fournisseur qui renvoie un enregistrement \emph{propre}
l'emporte. « Propre » veut dire : titre présent, au moins un
contributeur, et une année. Les correspondances partielles sont
enregistrées mais mybibli continue à parcourir la chaîne jusqu'à
trouver une correspondance propre ou épuiser les fournisseurs.

\subsection*{Zones bibliographiques : BnF et Library of Congress}

Six champs de catalogage détaillés — mention de responsabilité,
mention d'édition, titre et numéro de collection, note générale,
titre original — ne proviennent que des \emph{bibliothèques
nationales}. Les autres fournisseurs ne les portent pas.

Ces deux catalogues sont complémentaires plutôt que redondants : la
BnF décrit le fonds francophone, la Library of Congress le fonds
anglophone. Un titre absent des deux restera sans ces zones, et un
titre français que la BnF ne décrit pas ne sera pas rattrapé par la
Library of Congress.

Quand le fournisseur qui a résolu le titre ne remplit qu'une partie
de ces six zones, mybibli interroge l'autre bibliothèque pour
compléter les zones restées vides. La règle est \emph{la première
source l'emporte} : une valeur déjà présente n'est jamais écrasée.
Seules ces six zones sont fusionnées — le titre, les auteurs et la
couverture proviennent toujours d'un seul et même fournisseur.

Le bouton \emph{Compléter les métadonnées via les bibliothèques} de
l'onglet Santé applique le même traitement aux titres déjà
catalogués.

\subsection*{Replis sur les couvertures}

Quand le fournisseur qui résout le titre renvoie ses métadonnées
mais aucune URL de couverture (la BnF ne porte jamais d'URL
d'image dans son UNIMARC ; Google Books n'a parfois pas
d'\texttt{imageLinks} pour les titres auto-édités), mybibli essaie
trois sources dédiées aux couvertures, dans l'ordre, indexées par
ISBN :

\begin{enumerate}
  \item Le service \emph{Couvertures de la BnF}
        (\url{https://api.bnf.fr}) — les numérisations de
        couvertures du dépôt légal de la Bibliothèque nationale de
        France. La source la plus riche pour les éditeurs français
        et romands, avec une qualité d'image généralement
        excellente. Les images sont réutilisées selon la licence
        des métadonnées BnF (mention de la source).
  \item \emph{Inventaire.io} (\url{https://inventaire.io}) — une
        base de livres libre et communautaire ; particulièrement
        fournie en BD et manga francophones.
  \item L'API \emph{Open Library Covers}
        (\url{https://covers.openlibrary.org/}) — large couverture
        anglophone / multilingue.
\end{enumerate}

La première source qui possède une image l'emporte ; les replis ne
s'enclenchent que lorsque (a) la chaîne a réussi \emph{et} (b) le
fournisseur résolvant n'a pas fourni de \texttt{cover\_url}, de
sorte qu'un titre dont aucune source n'a la couverture finit sans
image, sans rafale d'essais inutiles.

\subsection*{Quand aucun fournisseur n'a la couverture : téléversez la vôtre}

Certains titres ne recevront jamais de couverture par la chaîne
automatique. En pratique, il s'agit surtout d'éditeurs français,
suisses et allemands, grand public comme universitaires (Dunod,
Eyrolles, Slatkine, Mardaga, Kana, Casterman, etc.) : Google Books
ne renvoie rien pour leurs ISBN et l'index Open Library Covers ne
les contient pas davantage. C'est une lacune structurelle des
sources de données gratuites, pas un bogue --- aucune récupération
répétée ne la comblera.

Pour ces titres, \textbf{le téléversement manuel est la voie
réaliste}. Sur la page de détail d'un titre, un bibliothécaire (ou
administrateur) voit un bouton \emph{Téléverser une couverture} sous
la zone de couverture ; les titres sans couverture reçoivent un
appel à l'action plus visible et un rappel d'une ligne sur cette
lacune. Les formats acceptés sont JPG, PNG, WebP et GIF, jusqu'à
10\,Mio. L'image est redimensionnée et stockée localement, et une
couverture téléversée manuellement est protégée contre tout
écrasement par une récupération automatique ultérieure.

Pour trouver les titres qui ont encore besoin d'une couverture,
utilisez le filtre \emph{Titres sans couverture} sur la page
d'accueil (bibliothécaire et plus), ou suivez le lien \emph{Voir les
titres sans couverture} affiché sous l'onglet Santé de la page
d'administration après une récupération en lot.

\section{Clés d'API}

Trois fournisseurs requièrent une clé d'API :

\subsection*{Google Books}

Optionnelle mais relève significativement le budget de requêtes
quotidien. Sans clé, le fournisseur Google Books reste désactivé —
la chaîne fonctionne toujours (BnF \texttt{→} Open Library
n'exigent pas de clé), mais les ouvrages anglophones grand public
que BnF ne référence pas arriveront avec des métadonnées vides.

Le parcours dans la Google Cloud Console déroute souvent les
nouveaux venus parce que l'étape de création de la clé est nichée
à deux menus de profondeur. Suivez exactement ces étapes :

\begin{enumerate}
  \item Ouvrez \url{https://console.cloud.google.com/} et
        connectez-vous.
  \item Sélecteur de projet en haut (à côté du logo
        \emph{Google Cloud}) → \emph{Nouveau projet} → nommez-le
        librement (p.\,ex.\ \texttt{mybibli-metadata}) →
        \emph{Créer}. Attendez quelques secondes que le projet
        soit sélectionné.
  \item Ouvrez le menu hamburger (en haut à gauche, trois lignes
        horizontales) → \emph{APIs \& services} →
        \emph{Bibliothèque}. Cherchez \textbf{Books API}, cliquez
        dessus, puis \emph{Activer}.
  \item Toujours sous \emph{APIs \& services}, ouvrez
        \emph{Identifiants}. Cliquez \emph{+ Créer des
        identifiants} → \emph{Clé d'API}. Une fenêtre affiche
        une clé fraîchement générée — copiez-la maintenant.
  \item Recommandé : cliquez \emph{Restreindre la clé} → sous
        \emph{Restrictions d'API} choisissez \emph{Restreindre la
        clé} → cochez uniquement \emph{Books API}. Laissez
        \emph{Restrictions d'application} sur \emph{Aucune}
        (mybibli n'envoie pas d'en-tête Referer). Enregistrez.
  \item Dans mybibli : \texttt{/admin > System > Metadata
        Providers} → collez dans le champ \textbf{Google Books
        API key} → \emph{Enregistrer}. Le prochain scan la prend
        en compte — pas de redémarrage requis.
\end{enumerate}

Pour vérifier que la clé est active, scannez un ISBN anglophone
que BnF ne connaît pas (\emph{Effective Java}, 9780134685991, est
un bon canari). Les métadonnées doivent se peupler en quelques
secondes. Si elles restent vides, voyez \emph{Dépannage} en fin
de chapitre.

\subsection*{OMDb}

\begin{enumerate}
  \item Visitez \url{https://www.omdbapi.com/apikey.aspx}.
  \item Choisissez le palier gratuit (1000 requêtes/jour).
  \item Confirmez via le courriel d'activation reçu.
  \item Collez la clé dans mybibli.
\end{enumerate}

\subsection*{TMDb}

\begin{enumerate}
  \item Inscrivez-vous sur \url{https://www.themoviedb.org/}.
  \item Ouvrez \emph{Settings → API} et demandez une clé pour
        \emph{usage personnel / non-commercial}.
  \item Collez la clé dans mybibli.
\end{enumerate}

\section{Où vivent les clés}

Les clés transitent par deux couches :

\begin{enumerate}
  \item \textbf{Variables d'environnement}\index{variables d'env!clés
        fournisseurs} (\texttt{GOOGLE\_BOOKS\_API\_KEY},
        \texttt{OMDB\_API\_KEY}, \texttt{TMDB\_API\_KEY}). Elles sont
        migrées \emph{une fois} au boot dans la table
        \texttt{settings} — elles servent essentiellement
        d'amorçage au moment du déploiement.
  \item \textbf{La table \texttt{settings}}, exposée via
        \texttt{/admin > System > Metadata Providers}. C'est la
        source de vérité au moment de la requête. La rotation d'une
        clé depuis l'interface prend effet à la toute prochaine
        récupération de métadonnées.
\end{enumerate}

L'implication est asymétrique :

\begin{itemize}
  \item Si vous positionnez une clé uniquement en variable
        d'environnement, le premier boot la copie dans la base et
        ensuite c'est la base qui fait foi.
  \item Si vous positionnez une clé uniquement dans l'interface,
        la variable d'environnement n'est plus jamais consultée.
  \item Si vous positionnez les deux puis \emph{Effacez} la clé
        depuis l'interface, la valeur en base devient vide — mais
        au boot suivant, la valeur d'environnement y est
        re-migrée. Pour rendre l'effacement durable, retirez
        aussi la variable d'environnement de \texttt{.env} /
        \texttt{docker-compose.yml}.
\end{itemize}

\section{Santé des fournisseurs}

\texttt{/admin > Health} pingue chaque fournisseur enregistré toutes
les cinq minutes et rapporte son état sous forme de coche verte ou
de croix rouge. Une croix rouge ne signifie pas toujours que le
fournisseur est en panne — cela peut aussi vouloir dire que la clé
d'API est absente ou invalide. L'onglet Health inclut le timestamp
du dernier check et le code HTTP pour faciliter le diagnostic.

\section{Quotas et back-pressure}

mybibli applique un petit quota par fournisseur (quelques requêtes
par seconde). Cataloguer cent volumes à la suite ne déclenchera pas
de 429 distant. Si un fournisseur renvoie quand même 429, la
requête est loguée et le fournisseur suivant de la chaîne est
consulté — le flux de catalogage ne se bloque jamais sur un
fournisseur en quota.

\begin{warning}
Ne committez jamais une vraie clé d'API dans git. Le fichier
\texttt{.env} est dans \texttt{.gitignore} pour cette raison ;
\texttt{.env.example} ne porte que des placeholders. Si une clé
fuit, faites-la tourner immédiatement sur la console du
fournisseur — coller une nouvelle clé dans mybibli prend effet en
quelques secondes.
\end{warning}

\section{Dépannage : métadonnées vides}

Si un ISBN fraîchement scanné finit avec des champs titre /
auteur / description vides, \textbf{interrogez d'abord
l'application}. Ouvrez la fiche et pressez \emph{Re-télécharger
les métadonnées} : la recherche se fait pendant que vous attendez,
et si elle revient vide, le message affiché dessous dit désormais
ce qui s'est réellement passé, au lieu de signaler seulement un
échec. Il distingue trois situations que rien ne séparait
jusqu'ici :

\begin{itemize}
  \item \textbf{«~\emph{N} sources ont été interrogées et aucune
        ne connaît ce code~»} \string— l'absence franche. Aucune
        clé, aucun réessai, aucune attente n'y changera quoi que
        ce soit : l'ouvrage n'est décrit dans aucun des catalogues
        que mybibli consulte. Saisissez les champs à la main.
  \item \textbf{«~\emph{Source} n'a pas pu répondre à
        l'instant~»} \string— la source était occupée ou trop
        lente. Là, il vaut réellement la peine de represser le
        bouton dans quelques minutes.
  \item \textbf{«~\emph{Source} n'a pas été consultée : aucune
        clé d'accès n'est renseignée~»} \string— elle n'a jamais
        été interrogée. Un champ dans
        \emph{Administration \textrightarrow{} Système} règle la
        question une fois pour toutes.
\end{itemize}

Le message nomme les sources, et ces noms ne sont jamais traduits :
ce sont ceux des bibliothèques elles-mêmes.

Si vous préférez raisonner par vous-même, parcourez cette liste :

\begin{enumerate}
  \item \textbf{État des fournisseurs dans \texttt{/admin >
        Health}.} Si Google Books ou Open Library affichent une
        croix rouge, ces fournisseurs sont sautés à chaque
        récupération. Vérifiez le code HTTP du dernier check —
        \texttt{403} ou \texttt{401} signifie en général une clé
        d'API absente ou invalide, pas une panne du fournisseur.
  \item \textbf{Langue du titre.} BnF est le fournisseur primaire
        et privilégie le français. Pour les livres anglophones,
        la chaîne bascule sur Google Books — donc la clé (voir
        plus haut) est la pièce qui débloque la situation.
  \item \textbf{Format de l'ISBN.} mybibli normalise EAN-13 vs
        ISBN-10 en interne, mais certains fournisseurs n'indexent
        que l'un des deux. Un second scan quelques minutes plus
        tard frappe parfois une route de repli différente.
  \item \textbf{Édition manuelle.} Vous pouvez toujours saisir
        titre, genre, contributeurs etc.\ à la main depuis la
        page de détail du titre — la sauvegarde fonctionne même
        si aucun fournisseur n'a répondu. Le genre \texttt{Non
        classé} est la valeur par défaut pour un titre arrivé
        sans données de classification.
\end{enumerate}

Si les champs restent vides après avoir configuré les bonnes
clés et attendu quelques secondes, ouvrez une issue sur le dépôt
du projet avec l'ISBN qui échoue — les mainteneurs ajoutent des
mocks de fournisseurs pour les cas reproductibles.
